\chapter{2000年京皖卷（春理）}
\section{单选题}
\begin{problem}
复数 \(z_{1} = 3 + \i\)，\(z_{2} = 1 - \i\)，则 \(z = z_{1} \cdot z_{2}\) 在复平面内的对应点位于（\quad）

\choices
    {第一象限}
    {第二象限}
    {第三象限}
    {第四象限}
\end{problem}

\begin{answer}
D
\end{answer}

\begin{solution}
\(z=(3+\i)(1-\i)=3-3\i+\i-\i^{2}=4-2\i\). 因为实部为正、虚部为负，所以对应点位于第四象限，故选 D.
\end{solution}

\begin{problem}
设全集 \(I = \{a, b, c, d, e\}\)，集合 \(M = \{a, c, d\}\)，\(N = \{b, d, e\}\)，那么 \(M \cap N\) 是（\quad）

\choices
    {\(\varnothing\)}
    {\(\{d\}\)}
    {\(\{a, c\}\)}
    {\(\{b, e\}\)}
\end{problem}

\begin{answer}
B
\end{answer}

\begin{solution}
集合 \(M\) 与集合 \(N\) 的公共元素只有 \(d\)，所以 \(M\cap N=\{d\}\)，故选 B.
\end{solution}

\begin{problem}
已知双曲线 \(\frac{x^{2}}{b^{2}} - \frac{y^{2}}{a^{2}} = 1\) 的两条渐近线互相垂直，那么该双曲线的离心率是（\quad）

\choices
    {\(2\)}
    {\(\sqrt{3}\)}
    {\(\sqrt{2}\)}
    {\(\frac{3}{2}\)}
\end{problem}

\begin{answer}
C
\end{answer}

\begin{solution}
该双曲线的渐近线为 \(y=\pm\frac{a}{b}x\). 两条渐近线垂直，所以 \(-\frac{a^{2}}{b^{2}}=-1\)，从而 \(a=b\). 双曲线的焦半距满足 \(c^{2}=a^{2}+b^{2}=2b^{2}\)，因此离心率 \(e=\frac{c}{b}=\sqrt{2}\)，故选 C.
\end{solution}

\begin{problem}
曲线 \(xy=1\) 的参数方程是（\quad）

\choices
    {\(\begin{cases}
        x=t^{\frac{1}{2}},\\
        y=t^{-\frac{1}{2}}
    \end{cases}\)}
    {\(\begin{cases}
        x=\sin\alpha,\\
        y=\csc\alpha
    \end{cases}\)}
    {\(\begin{cases}
        x=\cos\alpha,\\
        y=\sec\alpha
    \end{cases}\)}
    {\(\begin{cases}
        x=\tan\alpha,\\
        y=\cot\alpha
    \end{cases}\)}
\end{problem}

\begin{answer}
D
\end{answer}

\begin{solution}
选项 A 只给出第一象限的一支，选项 B、C 中的一个三角函数值只能覆盖有限范围，均不能表示整条曲线. 选项 D 满足 \(xy=\tan\alpha\cot\alpha=1\)，且 \(\tan\alpha\) 能取遍所有非零实数，正好覆盖曲线 \(xy=1\)，故选 D.
\end{solution}

\begin{problem}
一个圆锥的底面直径和高都同一个球的直径相等，那么圆锥与球的体积之比是（\quad）

\choices
    {\(1:3\)}
    {\(2:3\)}
    {\(1:2\)}
    {\(2:9\)}
\end{problem}

\begin{answer}
C
\end{answer}

\begin{solution}
设球的直径为 \(d\)，则圆锥的底面半径为 \(\frac{d}{2}\)，高为 \(d\). 因而圆锥体积为 \(\frac{1}{3}\pi\left(\frac{d}{2}\right)^{2}d=\frac{\pi d^{3}}{12}\)，球体积为 \(\frac{4}{3}\pi\left(\frac{d}{2}\right)^{3}=\frac{\pi d^{3}}{6}\). 所以体积之比为 \(1:2\)，故选 C.
\end{solution}

\begin{problem}
直线 \(\theta=\alpha\) 和直线 \(\rho\sin\left(\theta-\alpha\right)=1\) 的位置关系是（\quad）

\choices
    {垂直}
    {平行}
    {相交但不垂直}
    {重合}
\end{problem}

\begin{answer}
B
\end{answer}

\begin{solution}
由极坐标关系 \(x=\rho\cos\theta\)，\(y=\rho\sin\theta\)，有 \(\rho\sin\left(\theta-\alpha\right)=y\cos\alpha-x\sin\alpha=1\). 该直线的法向量为 \((-\sin\alpha,\cos\alpha)\)，所以它的方向与 \((\cos\alpha,\sin\alpha)\) 平行，即与直线 \(\theta=\alpha\) 平行. 又因前者不经过原点，二者不重合，故选 B.
\end{solution}

\begin{problem}
函数 \(y=\lg|x|\)（\quad）

\choices
    {是偶函数，在区间 \((-\infty,0)\) 上单调递增}
    {是偶函数，在区间 \((-\infty,0)\) 上单调递减}
    {是奇函数，在区间 \((0,+\infty)\) 上单调递增}
    {是奇函数，在区间 \((0,+\infty)\) 上单调递减}
\end{problem}

\begin{answer}
B
\end{answer}

\begin{solution}
因为 \(|-x|=|x|\)，所以 \(y=\lg|x|\) 是偶函数. 在 \((-\infty,0)\) 上，随着 \(x\) 增大，\(|x|\) 减小，而对数函数递增，因此 \(\lg|x|\) 单调递减，故选 B.
\end{solution}

\begin{problem}
从单词“equation”选取 5 个不同的字母排成一排，含有“qu”（其中“qu”相连且顺序不变）的不同排列共有（\quad）

\choices
    {\(120\) 个}
    {\(480\) 个}
    {\(720\) 个}
    {\(840\) 个}
\end{problem}

\begin{answer}
B
\end{answer}

\begin{solution}
\(q,u\) 必须同时选取并组成一个整体“qu”. 从其余 \(6\) 个字母中选取另外 \(3\) 个，有 \(\mathrm C_{6}^{3}\) 种选法；把“qu”和这 \(3\) 个字母排列，有 \(4!\) 种排法. 因此共有 \(\mathrm C_{6}^{3}\cdot4!=20\cdot24=480\) 种，故选 B.
\end{solution}

\begin{problem}
椭圆短轴长是 2，长轴长是短轴的 2 倍，则椭圆中心到其准线距离是（\quad）

\choices
    {\(\frac{8}{5}\sqrt{5}\)}
    {\(\frac{4}{5}\sqrt{5}\)}
    {\(\frac{8}{3}\sqrt{3}\)}
    {\(\frac{4}{3}\sqrt{3}\)}
\end{problem}

\begin{answer}
D
\end{answer}

\begin{solution}
设椭圆的长半轴、短半轴分别为 \(a\)，\(b\)，则 \(2b=2\)，\(2a=4\)，所以 \(a=2\)，\(b=1\). 于是焦半距 \(c=\sqrt{a^{2}-b^{2}}=\sqrt{3}\)，离心率 \(e=\frac{c}{a}=\frac{\sqrt{3}}{2}\). 中心到准线的距离为 \(\frac{a}{e}=\frac{4}{\sqrt{3}}=\frac{4}{3}\sqrt{3}\)，故选 D.
\end{solution}

\begin{problem}
函数 \(y=\frac{1}{2+\sin x+\cos x}\) 的最大值是（\quad）

\choices
    {\(\frac{\sqrt{2}}{2}-1\)}
    {\(\frac{\sqrt{2}}{2}+1\)}
    {\(1-\frac{\sqrt{2}}{2}\)}
    {\(-1-\frac{\sqrt{2}}{2}\)}
\end{problem}

\begin{answer}
B
\end{answer}

\begin{solution}
因为 \(\sin x+\cos x\geq-\sqrt{2}\)，所以分母满足 \(2+\sin x+\cos x\geq2-\sqrt{2}>0\). 因此函数最大值在分母取最小时取得，且 \(y_{\max}=\frac{1}{2-\sqrt{2}}=1+\frac{\sqrt{2}}{2}\)，故选 B.
\end{solution}

\begin{problem}
设复数 \(z_{1}=2\sin\theta+\i\cos\theta\)（\(\frac{\pi}{4}<\theta<\frac{\pi}{2}\)）在复平面上对应向量 \(\overrightarrow{OZ_{1}}\)，将 \(\overrightarrow{OZ_{1}}\) 按顺时针方向旋转 \(\frac{3\pi}{4}\) 后得到向量 \(\overrightarrow{OZ_{2}}\)，\(\overrightarrow{OZ_{2}}\) 对应的复数为 \(z_{2}=r\left(\cos\varphi+\i\sin\varphi\right)\)，则 \(\tan\varphi=\)（\quad）

\choices
    {\(\frac{2\tan\theta+1}{2\tan\theta-1}\)}
    {\(\frac{2\tan\theta-1}{2\tan\theta+1}\)}
    {\(\frac{1}{2\tan\theta+1}\)}
    {\(\frac{1}{2\tan\theta-1}\)}
\end{problem}

\begin{answer}
A
\end{answer}

\begin{solution}
\(z_{1}\) 对应点的坐标为 \((2\sin\theta,\cos\theta)\). 顺时针旋转 \(\frac{3\pi}{4}\) 后，所得向量的坐标为 \(\left(\frac{\cos\theta-2\sin\theta}{\sqrt{2}},-\frac{2\sin\theta+\cos\theta}{\sqrt{2}}\right)\). 因而 \(\tan\varphi=\frac{-(2\sin\theta+\cos\theta)}{\cos\theta-2\sin\theta}=\frac{2\tan\theta+1}{2\tan\theta-1}\)，故选 A.
\end{solution}

\begin{problem}
设 \(\alpha\)，\(\beta\) 是一个钝角三角形的两个锐角，下列四个不等式中不正确的是（\quad）

\choices
    {\(\tan\alpha\tan\beta<1\)}
    {\(\sin\alpha+\sin\beta<\sqrt{2}\)}
    {\(\cos\alpha+\cos\beta>1\)}
    {\(\frac{1}{2}\tan\left(\alpha+\beta\right)<\tan\frac{\alpha+\beta}{2}\)}
\end{problem}

\begin{answer}
D
\end{answer}

\begin{solution}
设 \(S=\alpha+\beta\). 因为第三个角为钝角，所以 \(0<S<\frac{\pi}{2}\). 由正切加法公式
\(\tan S=\frac{\tan\alpha+\tan\beta}{1-\tan\alpha\tan\beta}>0\)，且 \(\tan\alpha+\tan\beta>0\)，所以 \(1-\tan\alpha\tan\beta>0\)，故 \(\tan\alpha\tan\beta<1\). 又因为 \(\lvert\alpha-\beta\rvert<S<\frac{\pi}{2}\)，所以
\(\sin\alpha+\sin\beta=2\sin\frac{S}{2}\cos\frac{\alpha-\beta}{2}<2\sin\frac{\pi}{4}=\sqrt{2}\).
同时 \(0\leq\frac{\lvert\alpha-\beta\rvert}{2}<\frac{S}{2}<\frac{\pi}{4}\)，由余弦函数在 \([0,\frac{\pi}{2}]\) 上单调递减，得
\(\cos\frac{\alpha-\beta}{2}>\cos\frac{S}{2}\)，从而 \(\cos\alpha+\cos\beta=2\cos\frac{S}{2}\cos\frac{\alpha-\beta}{2}>2\cos^{2}\frac{S}{2}=1+\cos S>1\). 最后令 \(t=\tan\frac{S}{2}\)，则 \(0<t<1\)，所以 \(\frac{1}{2}\tan S=\frac{t}{1-t^{2}}>t=\tan\frac{S}{2}\). 因此选项 D 的不等式不成立，故选 D.
\end{solution}

\begin{problem}
已知等差数列 \(\{a_{n}\}\) 满足 \(a_{1}+a_{2}+a_{3}+\dots+a_{101}=0\)，则有（\quad）

\choices
    {\(a_{1}+a_{101}>0\)}
    {\(a_{2}+a_{100}<0\)}
    {\(a_{3}+a_{99}=0\)}
    {\(a_{51}=51\)}
\end{problem}

\begin{answer}
C
\end{answer}

\begin{solution}
等差数列的前 \(101\) 项和为 \(101a_{51}\)，所以 \(a_{51}=0\). 对称项满足 \(a_k+a_{102-k}=2a_{51}=0\)，特别地 \(a_{3}+a_{99}=0\). 因此正确选项为 C.
\end{solution}

\begin{problem}
已知函数 \(f(x)=ax^{3}+bx^{2}+cx+d\) 的图象如图，则（\quad）

\bitmapfigure[max width=0.45\linewidth]{img/2000/jing_hui_science_spring/q14_fig1.png}

\choices
    {\(b\in(-\infty,0)\)}
    {\(b\in(0,1)\)}
    {\(b\in(1,2)\)}
    {\(b\in(2,+\infty)\)}
\end{problem}

\begin{answer}
A
\end{answer}

\begin{solution}
由图象可知函数的三个零点为 \(0\)，\(1\)，\(2\)，且图象右端上升，所以 \(a>0\). 因而 \(f(x)=a x(x-1)(x-2)=a(x^{3}-3x^{2}+2x)\)，从而 \(b=-3a<0\)，故选 A.
\end{solution}

\section{填空题}
\begin{problem}
函数 \(y=\cos\left(\frac{2\pi}{3}x+\frac{\pi}{4}\right)\) 的最小正周期是\fillinblank{}.
\end{problem}

\begin{answer}
\(3\)
\end{answer}

\begin{solution}
余弦函数 \(\cos(\omega x+\varphi)\) 的最小正周期为 \(\frac{2\pi}{|\omega|}\). 此题 \(|\omega|=\frac{2\pi}{3}\)，所以最小正周期为 \(\frac{2\pi}{2\pi/3}=3\).
\end{solution}

\begin{problem}
已知一体积为 \(72\) 的正四面体，连接两个面的重心 \(E\)，\(F\)，则线段 \(EF\) 的长是\fillinblank{}.
\end{problem}

\begin{answer}
\(2\sqrt{2}\)
\end{answer}

\begin{solution}
设正四面体的棱长为 \(s\). 每个面的面积为 \(\frac{\sqrt{3}}{4}s^{2}\)，正四面体的高为 \(s\sqrt{\frac{2}{3}}\)，所以 \(72=\frac{1}{3}\cdot\frac{\sqrt{3}}{4}s^{2}\cdot s\sqrt{\frac{2}{3}}=\frac{s^{3}}{6\sqrt{2}}\)，从而 \(s=6\sqrt{2}\). 若两个面分别为 \(ABC\)，\(ABD\)，其重心为 \(E\)，\(F\)，则 \(\overrightarrow{EF}=\frac{1}{3}\overrightarrow{CD}\)，所以 \(EF=\frac{s}{3}=2\sqrt{2}\).
\end{solution}

\begin{problem}
\(\left(\sqrt{x}-\frac{1}{\sqrt[3]{x}}\right)^{10}\) 展开式中的常数项是\fillinblank{}.
\end{problem}

\begin{answer}
\(210\)
\end{answer}

\begin{solution}
通项取第二项 \(k\) 次时为 \((-1)^{k}\mathrm C_{10}^{k}x^{5-5k/6}\). 要成为常数项，必须有 \(5-\frac{5k}{6}=0\)，所以 \(k=6\). 因此常数项为 \((-1)^{6}\mathrm C_{10}^{6}=210\).
\end{solution}

\begin{problem}
在空间，有下列命题：

\begin{circlelist}
\item 如果两直线 \(a\)，\(b\) 分别与直线 \(l\) 平行，那么 \(a\parallel b\)；
\item 如果直线 \(a\) 与平面 \(\beta\) 内的一条直线 \(b\) 平行，那么 \(a\parallel\beta\)；
\item 如果直线 \(a\) 与平面 \(\beta\) 内的两条直线 \(b\)，\(c\) 都有垂直，那么 \(a\perp\beta\)；
\item 如果平面 \(\beta\) 内的一条直线 \(a\) 垂直平面 \(\gamma\)，那么 \(\beta\perp\gamma\).
\end{circlelist}

其中正确的命题的序号是\fillinblank{}.
\end{problem}

\begin{answer}
\(\circled{1}\)，\(\circled{4}\)
\end{answer}

\begin{solution}
命题 \(\circled{1}\) 中，两条直线都与同一直线平行，因而方向相同且共面，所以 \(a\parallel b\)，正确. 命题 \(\circled{2}\) 中，直线 \(a\) 也可能在平面 \(\beta\) 内，不能据此断定 \(a\parallel\beta\)，错误. 命题 \(\circled{3}\) 还缺少直线 \(b\)，\(c\) 相交的条件，故不能断定 \(a\perp\beta\)，错误. 命题 \(\circled{4}\) 中，平面 \(\beta\) 含有一条垂直于平面 \(\gamma\) 的直线，所以 \(\beta\perp\gamma\)，正确. 因此填 \(\circled{1}\)，\(\circled{4}\).
\end{solution}

\section{解答题}
\begin{problem}
在 \(\triangle ABC\) 中，角 \(A\)，\(B\)，\(C\) 对边分别为 \(a\)，\(b\)，\(c\)，证明：\(\frac{a^{2}-b^{2}}{c^{2}}=\frac{\sin\left(A-B\right)}{\sin C}\).
\end{problem}

\begin{solution}
由正弦定理，存在外接圆半径 \(R\)，使得 \(a=2R\sin A\)，\(b=2R\sin B\)，\(c=2R\sin C\). 因而 \(\frac{a^{2}-b^{2}}{c^{2}}=\frac{\sin^{2}A-\sin^{2}B}{\sin^{2}C}\).

利用恒等式 \(\sin^{2}A-\sin^{2}B=\sin(A+B)\sin(A-B)\)，以及 \(A+B=\pi-C\)，可得 \(\frac{a^{2}-b^{2}}{c^{2}}=\frac{\sin(A+B)\sin(A-B)}{\sin^{2}C}=\frac{\sin C\sin(A-B)}{\sin^{2}C}=\frac{\sin(A-B)}{\sin C}\).
因此等式成立.
\end{solution}

\begin{problem}
在直角梯形 \(ABCD\) 中，\(\angle D=\angle BAD=90^{\circ}\)，\(AD=DC=\frac{1}{2}AB=a\)（如图一）. 将 \(\triangle ADC\) 沿 \(AC\) 折起，使 \(D\) 到 \(D'\). 记面 \(ACD'\) 为 \(\alpha\)，面 \(ABC\) 为 \(\beta\)，面 \(BCD'\) 为 \(\gamma\).

\begin{enumerate}
\item 若二面角 \(\alpha-AC-\beta\) 为直二面角（如图二），求二面角 \(\beta-BC-\gamma\) 的大小；
\item 若二面角 \(\alpha-AC-\beta\) 为 \(60^{\circ}\)（如图三），求三棱锥 \(D'-ABC\) 的体积.
\end{enumerate}

\FigureLayoutDeclare{img/2000/jing_hui_science_spring/q20_fig1.png}{12.0cm}{40bp 65bp 33bp 38bp}
\bitmapfigure[max width=0.9\linewidth]{img/2000/jing_hui_science_spring/q20_fig1.png}
\end{problem}

\begin{solution}
\begin{enumerate}
\item 取直角坐标系使 \(A=(0,0,0)\)，\(B=(2a,0,0)\)，\(C=(a,a,0)\)，\(D=(0,a,0)\). 设 \(H\) 为 \(D\) 在 \(AC\) 上的射影，则 \(H=\left(\frac{a}{2},\frac{a}{2},0\right)\)，且 \(DH=\frac{a}{\sqrt{2}}\). 当二面角 \(\alpha-AC-\beta\) 为直二面角时，绕轴 \(AC\) 将 \(HD\) 旋转 \(90^{\circ}\)，可取 \(D'=\left(\frac{a}{2},\frac{a}{2},\frac{a}{\sqrt{2}}\right)\).

取平面 \(\beta\) 和平面 \(\gamma\) 的法向量，分别可取 \(\bs n_{\beta}=(0,0,1)\) 及 \(\bs n_{\gamma}=\overrightarrow{BC}\times\overrightarrow{BD'}=a^{2}\left(\frac{1}{\sqrt{2}},\frac{1}{\sqrt{2}},1\right)\). 因此二面角 \(\beta-BC-\gamma\) 的锐角 \(\theta\) 满足 \(\cos\theta=\frac{|\bs n_{\beta}\cdot\bs n_{\gamma}|}{|\bs n_{\beta}|\,|\bs n_{\gamma}|}=\frac{1}{\sqrt{2}}\)，所以 \(\theta=45^{\circ}\).

\item 三角形 \(ABC\) 的面积为 \(\frac{1}{2}\cdot 2a\cdot a=a^{2}\). 折起角为 \(60^{\circ}\) 时，点 \(D'\) 到平面 \(ABC\) 的距离为 \(DH\sin60^{\circ}=\frac{a}{\sqrt{2}}\cdot\frac{\sqrt{3}}{2}=\frac{a\sqrt{6}}{4}\).

所以三棱锥 \(D'-ABC\) 的体积为 \(V=\frac{1}{3}S_{\triangle ABC}\cdot\frac{a\sqrt{6}}{4}=\frac{\sqrt{6}}{12}a^{3}\).
\end{enumerate}
\end{solution}

\begin{problem}
设函数 \(f(x)=|\lg x|\)，若 \(0<a<b\)，且 \(f(a)>f(b)\)，证明：\(ab<1\).
\end{problem}

\begin{solution}
若 \(b\leq1\)，则 \(0<a<b\leq1\)，从而 \(ab<b^{2}\leq1\)，且因 \(a<b\) 有 \(ab<1\). 若 \(b>1\)，则不能有 \(a\geq1\)，因为这时 \(f(x)=\lg x\) 在 \([1,+\infty)\) 上递增，会得到 \(f(a)<f(b)\). 因此 \(0<a<1<b\)，于是 \(f(a)=-\lg a\)，\(f(b)=\lg b\). 由题设 \(-\lg a>\lg b\)，所以 \(\lg(ab)<0\)，即 \(ab<1\).
\end{solution}

\begin{problem}
如图，设点 \(A\) 和 \(B\) 为抛物线 \(y^{2}=4px\)（\(p>0\)）上原点以外的两个动点，已知 \(OA\perp OB\)，\(OM\perp AB\). 求点 \(M\) 的轨迹方程，并说明它表示什么曲线.

\bitmapfigure[max width=0.45\linewidth]{img/2000/jing_hui_science_spring/q22_fig1.png}
\end{problem}

\begin{solution}
设抛物线上点的参数分别为 \(t\)，\(u\)，则 \(A=(pt^{2},2pt)\)，\(B=(pu^{2},2pu)\)，且 \(t\neq0\)，\(u\neq0\). 由 \(OA\perp OB\) 得 \(\overrightarrow{OA}\cdot\overrightarrow{OB}=p^{2}tu(tu+4)=0\)，所以 \(tu=-4\).

参数为 \(t\)，\(u\) 的抛物线弦 \(AB\) 的方程为 \(2x-(t+u)y+2ptu=0\)，代入 \(tu=-4\) 得 \(2x-(t+u)y-8p=0\). 令 \(L=t+u\)，则直线 \(AB\) 为 \(2x-Ly-8p=0\). 原点到这条直线的垂足为 \(M=\left(\frac{16p}{L^{2}+4},-\frac{8pL}{L^{2}+4}\right)\).

于是 \(x^{2}+y^{2}=\frac{64p^{2}}{L^{2}+4}=4px\)，所以轨迹满足 \(x^{2}+y^{2}-4px=0\). 反之，任意实数 \(L\) 都能取为方程 \(v^{2}-Lv-4=0\) 的两根之和，且两根乘积为 \(-4\)，故上述参数能得到圆上除原点外的所有点. 因而轨迹是方程 \(x^{2}+y^{2}-4px=0\) 所表示的圆去掉原点，即圆心为 \((2p,0)\)、半径为 \(2p\) 的圆去掉原点.
\end{solution}

\begin{problem}
某地区上年度电价为 \(0.8\text{ 元}/(\mathrm{kW}\cdot\mathrm{h})\)，年用电量为 \(a\mathrm{~kW\cdot h}\). 本年度计划将电价降到 \(0.55\text{ 元}/(\mathrm{kW}\cdot\mathrm{h})\) 至 \(0.75\text{ 元}/(\mathrm{kW}\cdot\mathrm{h})\) 之间，而用户期望电价为 \(0.4\text{ 元}/(\mathrm{kW}\cdot\mathrm{h})\). 经测算，下调电价后新增的用电量与实际电价和用户期望电价的差成反比（比例系数为 \(k\)）. 该地区电力的成本为 \(0.3\text{ 元}/(\mathrm{kW}\cdot\mathrm{h})\).

\begin{enumerate}
\item 写出本年度电价下调后，电力部门的收益 \(y\) 与实际电价 \(x\) 的函数关系式；
\item 设 \(k=0.2a\)，当电价最低定为多少时，仍可保证电力部门的收益比上年至少增长 20\%.
\end{enumerate}

注：收益=实际用电量 \(\times\)（实际电价-成本价）.
\end{problem}

\begin{solution}
\begin{enumerate}
\item 设本年度实际电价为 \(x\)，则 \(0.55\leq x\leq0.75\). 因为新增用电量与 \(x-0.4\) 成反比，且比例系数为 \(k\)，所以新增用电量为 \(\frac{k}{x-0.4}\)，本年度实际用电量为 \(a+\frac{k}{x-0.4}\). 每千瓦时的收益为 \(x-0.3\)，故 \(y=\left(a+\frac{k}{x-0.4}\right)(x-0.3)\)，其中 \(0.55\leq x\leq0.75\).

\item 上年度收益为 \(a(0.8-0.3)=0.5a\)，本年度收益至少增长 \(20\%\)，所以必须有 \(y\geq1.2\times0.5a=0.6a\). 令 \(k=0.2a\)，并设 \(u=x-0.4\)，则 \(u\in[0.15,0.35]\)，且 \(\frac{y}{a}=\left(1+\frac{0.2}{u}\right)(u+0.1)=u+0.3+\frac{0.02}{u}\).

由 \(y\geq0.6a\) 且 \(u>0\)，得 \(u+0.3+\frac{0.02}{u}\geq0.6\)，即 \(u^{2}-0.3u+0.02\geq0\)，也就是 \((u-0.1)(u-0.2)\geq0\). 在 \([0.15,0.35]\) 内只能有 \(u\geq0.2\)，所以 \(x\geq0.6\). 因此电价最低定为 \(0.6\text{ 元}/(\mathrm{kW}\cdot\mathrm{h})\).
\end{enumerate}
\end{solution}

\begin{problem}
已知函数 \(f(x)=\begin{cases}
f_{1}(x),&x\in\left[0,\frac{1}{2}\right),\\
f_{2}(x),&x\in\left[\frac{1}{2},1\right],
\end{cases}\)
其中 \(f_{1}(x)=-2\left(x-\frac{1}{2}\right)^{2}+1\)，\(f_{2}(x)=-2x+2\).

\begin{enumerate}
\item 在下面坐标系上画出 \(y=f(x)\) 的图象；
\item 设 \(y=f_{2}(x)\)（\(x\in\left[\frac{1}{2},1\right]\)）的反函数为 \(y=g(x)\)，\(a_{1}=1\)，\(a_{2}=g(a_{1})\)，\(\dots\)，\(a_{n}=g(a_{n-1})\)，求数列 \(\{a_{n}\}\) 的通项公式，并求 \(\displaystyle\lim_{n\to\infty}a_{n}\)；
\item 若 \(x_{0}\in\left[0,\frac{1}{2}\right)\)，\(x_{1}=f(x_{0})\)，\(f(x_{1})=x_{0}\)，求 \(x_{0}\).
\end{enumerate}

\bitmapfigure[max width=0.45\linewidth]{img/2000/jing_hui_science_spring/q24_fig1.png}
\end{problem}

\begin{solution}
\begin{enumerate}
\item 当 \(0\leq x<\frac{1}{2}\) 时，\(f(x)=f_{1}(x)\)，且 \(f_{1}'(x)=2-4x>0\)，所以图象是抛物线 \(y=-2\left(x-\frac{1}{2}\right)^{2}+1\) 从实心点 \(\left(0,\frac{1}{2}\right)\) 出发的相应部分，单调上升并趋近于开点 \(\left(\frac{1}{2},1\right)\). 当 \(\frac{1}{2}\leq x\leq1\) 时，\(f(x)=f_{2}(x)=-2x+2\)，图象是从实心点 \(\left(\frac{1}{2},1\right)\) 到实心点 \((1,0)\) 的线段. 因此图象由上述抛物线段和线段组成，在 \(\left(\frac{1}{2},1\right)\) 处取线段的实心点.

\item 由 \(y=-2x+2\)，得反函数 \(g(x)=1-\frac{x}{2}\)，其定义域为 \([0,1]\). 数列满足 \(a_n=1-\frac{a_{n-1}}{2}\quad(n\geq2)\)，且 \(a_1=1\). 设该递推的平衡值为 \(L\)，则 \(L=1-\frac{L}{2}\)，所以 \(L=\frac{2}{3}\). 因而 \(a_n-L=-\frac{1}{2}(a_{n-1}-L)\)，从而 \(a_n-\frac{2}{3}=\left(-\frac{1}{2}\right)^{n-1}\left(1-\frac{2}{3}\right)\). 所以 \(a_n=\frac{2}{3}+\frac{1}{3}\left(-\frac{1}{2}\right)^{n-1}\)，并且 \(\displaystyle\lim_{n\to\infty}a_n=\frac{2}{3}\).

\item 因为 \(x_0\in\left[0,\frac{1}{2}\right)\)，所以 \(x_1=f(x_0)=f_1(x_0)\in\left[\frac{1}{2},1\right)\)，从而 \(f(x_1)=f_2(x_1)\). 代入得 \(x_0=-2\left[-2\left(x_0-\frac{1}{2}\right)^2+1\right]+2=4\left(x_0-\frac{1}{2}\right)^2\).
于是 \(4x_0^2-5x_0+1=0\)，即 \((4x_0-1)(x_0-1)=0\). 结合 \(x_0\in[0,\frac{1}{2})\)，得到 \(x_0=\frac{1}{4}\). 代回可得 \(x_1=\frac{7}{8}\)，且 \(f(x_1)=\frac{1}{4}\)，满足条件.
\end{enumerate}
\end{solution}
